\ifdefined\mainfile\else
\documentclass{article}
\usepackage{graphicx}
\usepackage{float}
\usepackage{amsmath}
\usepackage{amssymb}
\usepackage[a4paper, left=2.5cm, right=2.5cm, top=2.5cm, bottom=2.5cm]{geometry}
\usepackage[colorlinks=true,linkcolor=black,citecolor=blue,urlcolor=blue]{hyperref}
\begin{document}
\fi

\section{Instruction Decoder / Controller}
\label{sec:idc}
% State machine med case-in-case struktur
% States: INF, EX0, EX1, EX2, EX3, EX4
% Genererer kontrolord: PS, IL, DX, AX, BX, MB, FS, MD, RW, MM, MW
% Modtager status flags: V, C, N, Z

% \begin{figure}[H]
%     \centering
%     \includegraphics[width=0.75\linewidth]{images/idc/TODO_state_diagram.png}
%     \caption{IDC state diagram}
%     \label{fig:idc-states}
% \end{figure}

\subsection{State Machine Design}
\label{sec:idc-fsm}
% State diagram, state transitions, opcode dekodning
% Case-in-case: ydre case = current_state, indre case = opcode

\subsection{Kontrolord}
\label{sec:idc-controlword}
% 28-bit kontrolord og hvordan det genereres pr. state/opcode kombination
% Tabel over kontrolord for hver instruktion

\subsection{VHDL Implementering}
\label{sec:idc-vhdl}
% VHDL kode og forklaring af to-process model (state register + next state/output logic)

\subsection{Test}
\label{sec:idc-test}
% Testbench og verifikation -- ADD, LD, BRZ, JMP, etc.

\ifdefined\mainfile\else
\end{document}
\fi
